\documentclass[a4paper, titlepage]{book}

\usepackage{../preamble}


\begin{document}
\frontmatter

\titlePage{2}{Electricity and Magnetism I}

\tableofcontents


\mainmatter

\chapter{Electric Charges, Fields, and Gauss's Law}

\section{Electric Charges}

\textbf{Electric charge} is denoted by the symbol \(q\) or sometimes \(Q\),
and its SI unit is the coulomb (C).

Charge can be either positive or negative,
with protons carrying a positive charge and electrons carrying a negative charge.

The \textbf{law of conservation of  charge} states that charge can not created or destroyed,
only transferred,
and the net charge has to remain constant.
There are two ways to charge and discharge (i.e., move charges):
\begin{itemize}
    \item Conduction -- friction (triboelectricity) or contact
    \item Induction -- contactless
\end{itemize}

\section{Coulomb's Law}

Charged objects exert an \textbf{electric force} on other charged objects around them.
The magnitude of the electric force can be given by \textbf{Coulomb's law}:
\begin{equation} \label{eq:coulombs-law}
    F = k\frac{q_1 q_2}{r^2}
\end{equation}
where
\begin{itemize}
    \item \(F\) is the force \emph{both} objects experience as given by Newton's Third law.
    \item \(k \approx \qty{8.99e-9}{\newton \meter\squared \per \coulomb\squared}\)
          and is called the Coulomb constant.
    \item \(q_1\) and \(q_2\) are the two charges of the two objects.
          Note that both values are signed.
    \item \(r\) is the distance between the two particles.
\end{itemize}
Objects with charges of the same sign will repel each other,
while objects with charges of opposite signs will attract each other.
Thus, Coulomb's law will give a \emph{negative} value for attraction
and a \emph{positive} value for repulsion.

\section{Electric Flux}

\textbf{Flux}, symbolized by \(\Phi\),
is the amount of field that passes through a \textbf{surface},
which may be imaginary.
The general formula for the electric flux \(\Phi_E\) is
\begin{equation}
    \Phi_E = \iint_S \bm{E}(\bm{r}) \cdot \odif{\bm{S}}
\end{equation}
where
\begin{itemize}
    \item \(S\) is the surface to integrate over.
    \item \(\bm{E}(\bm{r})\) is the function of electric field in a 3D space.
\end{itemize}

For introductory electricity and magnetism,
we usually focus on a few simpler cases.

Surfaces can be either open (e.g., a window pane) or closed (e.g., a box).
For the case of a constant electric field passing through an flat open surface,
the flux equation simplifies to
\begin{equation}
    \begin{aligned}
        \Phi = \bm{E} \cdot \bm{A} \\
        \Phi = E A \cos\theta
    \end{aligned}
\end{equation}
where
\begin{itemize}
    \item \(\bm{E}\) is the electric field vector, and \(E\) is its magnitude.
    \item \(\bm{A}\) is area vector, and \(A\) is its magnitude.
\end{itemize}


\chapter{Electric Potential}

Electric potential energy is denoted by \(U_\mathrm{E}\)


\chapter{Conductors and Capacitors}


\chapter{Electric Circuits}

\section{Ohm's Law}

\section{Gauss's Law of Magnetism}


\chapter{Magnetic Fields and Electromagnetism}

\section{Ampère's Law}

\section{Biot-Savart Law}


\chapter{Electromagnetic Induction}

\section{Faraday's Law and Lenz's Law}

\section{Maxwell's Equations}


\appendix


\chapter{Vectors} \label{ch:vectors}

\begin{center}
    \begin{tabular}{ c|c }
    \textbf{Vector}                                    & \textbf{Scalar} \\
    \hline
    Has both \textbf{direction} and \textbf{magnitude} & Has only magnitude \\
    Displacement                                       & Length/distance \\
    Velocity                                           & Speed \\
    Gravitational fields                               & Mass
    \end{tabular}
\end{center}

\section{Notation}

\begin{subequations}
    2D:
    \begin{equation}
        \vec{v}
        = \bm{v}
        = \ab< v_x, v_y >
        = \ab( v_x, v_y )
        = v_x \bm{\hat{\imath}} + v_y \bm{\hat{\jmath}}
    \end{equation}
    3D:
    \begin{equation}
        \vec{v}
        = \bm{v}
        = \ab< v_x, v_y, v_z >
        = \ab( v_x, v_y, v_z )
        = v_x \bm{\hat{\imath}} + v_y \bm{\hat{\jmath}} + v_z \bm{\hat{k}}
    \end{equation}
\end{subequations}

Note that \(\vec{v}\) is more common in handwriting,
and \(\bm{v}\) is more common in print.

See \pgRef{sec:unit-vectors} for information about \(\bm{\hat{\imath}}\),
\(\bm{\hat{\jmath}}\), and \(\bm{\hat{k}}\).

\section{Magnitude}

\textbf{Magnitude} is the size of the vector.

\begin{subequations}
    2D:
    \begin{equation}
        \ab\| \bm{v} \| = \sqrt{{v_x}^2 + {v_y}^2}
    \end{equation}
    3D:
    \begin{equation}
        \ab\| \bm{v} \| = \sqrt{{v_x}^2 + {v_y}^2 + {v_z}^2}
    \end{equation}
\end{subequations}

\section{Unit Vectors} \label{sec:unit-vectors}

\textbf{Units vectors} are vectors with a magnitude of 1.
They are denoted using a "hat", e.g., \(\hat{v}\) is "v-hat".

The \textbf{standard unit vectors} are
\begin{equation}
    \begin{aligned}
        \bm{\hat{\imath}} &= \bm{\hat{x}} = \ab< 1, 0, 0 >, \\
        \bm{\hat{\jmath}} &= \bm{\hat{y}} = \ab< 0, 1, 0 >, \\
        \bm{\hat{k}} &= \bm{\hat{z}} = \ab< 0, 0, 1 >
    \end{aligned}
\end{equation}
\vfil

\section{Operations}

Vector addition:
\begin{equation}
    \bm{u} + \bm{v} = \ab< u_x + v_x, u_y + v_y, u_z + v_z >
\end{equation}

Scalar multiplication:
\begin{equation}
    c \bm{v} = \ab< c v_x, c v_y, c v_z >
\end{equation}

Dot product:
\begin{equation}
    \begin{aligned}
        \bm{u} \cdot \bm{v} &= u_x v_x + u_y v_y + u_z v_z \\
        &= \ab\| \bm{u} \| \ab\| \bm{v} \| \cos\theta
    \end{aligned}
\end{equation}

Cross product:
\begin{equation}
    \begin{aligned}
        \bm{u} \times \bm{v}
        &= \ab< u_y v_z - u_z v_y, u_z v_a - u_a v_z, u_a v_y - u_y v_a > \\
        &= \ab\| \bm{u} \| \ab\| \bm{v} \| \bm{\hat{w}} \sin\theta \\
        &= \bm{w}
    \end{aligned}
\end{equation}
where \(\bm{\hat{w}}\) is a unit vector orthogonal to both \(\bm{u}\) and \(\bm{v}\)
as given by the \textbf{right-hand rule}.
Hold your index, middle, and thumb fingers such that all three are perpendicular to each other.
Your index finger represents \(\bm{u}\), your middle finger \(\bm{v}\),
and your thumb \(\bm{w}\).



\backmatter

\printbibliography[
    heading=bibintoc,
    title={References}
]


\end{document}